% Generated draft: requires mathematical and source review.
\documentclass[11pt]{article}
\usepackage[margin=25mm]{geometry}
\usepackage{fontspec}
\setmainfont{Latin Modern Roman}
\usepackage{amsmath,amssymb,mathtools,unicode-math}
\setmathfont{Latin Modern Math}
\usepackage{enumitem,needspace,xurl,hyperref}
\setlength{\parindent}{0pt}
\setlength{\parskip}{5pt}
\setlength{\emergencystretch}{3em}
\newcommand{\N}{\mathbb{N}}
\newcommand{\Z}{\mathbb{Z}}
\newcommand{\Q}{\mathbb{Q}}
\newcommand{\R}{\mathbb{R}}
\newcommand{\C}{\mathbb{C}}
\newcommand{\F}{\mathbb{F}}
\newcommand{\A}{\mathbb{A}}
\newcommand{\PP}{\mathbb{P}}
\newcommand{\E}{\mathbb{E}}
\newcommand{\eps}{\varepsilon}
\newcommand{\dd}{\,\mathrm{d}}
\begin{document}
\section*{UnsolvedMath ID: 30004666}
\textbf{Source problem code:} catalogue entry
\subsection{Definitions and mathematical statement}
Let \(F\) be a Laurent-series field of characteristic \(p>0\), and let \(G\) be an absolutely simple algebraic group over \(F\). Interpret \(\pi_1(G)\) as the algebraic fundamental group: after passage to a separable closure and choice of a maximal torus \(T\), it is the quotient of the cocharacter lattice \(X_*(T)\) by the coroot lattice. Assume
\[
p\mid |\pi_1(G)|.
\]
Let \(a\) be an alcove in the relevant Bruhat--Tits apartment, let \(\operatorname{Fl}_{G,a}\) denote the full affine flag variety associated with its Iwahori subgroup, and let \(\operatorname{Fl}_{G,a}^{\circ}\) denote its neutral component.

In a setting where the Schubert varieties in this component are indexed by the affine Weyl group \(W_{\mathrm{aff}}\), write \(C_w\) for the Iwahori orbit indexed by \(w\), and give its closure the reduced induced scheme structure:
\[
X_w=\overline{C_w}_{\mathrm{red}}\subseteq\operatorname{Fl}_{G,a}^{\circ}.
\]
A scheme is normal if each of its local rings is an integrally closed domain. The intended classification problem in the supplied corrected statement is to determine, for every intended pair \((G,p)\), the subset
\[
\mathcal N_{G,a}
=\{w\in W_{\mathrm{aff}}:X_w\text{ is normal}\}.
\]
Equivalently, using the base alcove \(a\), classify the alcoves \(w(a)\) whose corresponding Schubert varieties are normal.

\textit{Unresolved scope qualifications.} The record does not specify the coefficient or residue field, splitting or ramification hypotheses on \(G\), or the full list referred to as the ``relevant pairs \((G,p)\)'' in the original excerpt. It also does not specify whether normality is intended over the residue field or after extension to an algebraic closure, nor whether the orbit indexing is formulated over that field or geometrically. The displayed classification therefore states the intended target conditional on specifying this affine-flag setup; it does not supply these missing hypotheses. The original excerpt asks to classify ``all Schubert varieties'' without the word ``normal''; the corrected statement and title explicitly identify normal Schubert varieties as the intended target.
\subsection{Short English statement}
For each intended absolutely simple group \(G\) over a Laurent-series field of characteristic \(p>0\) with \(p\mid|\pi_1(G)|\), classify the normal Schubert varieties in the neutral component of its full affine flag variety, or their affine Weyl-group elements and alcoves. The intended pairs, field hypotheses, and geometric interpretation require clarification. The supplied record labels the problem partially solved; this draft does not independently verify that assessment.
\subsection{Sources}
\begin{itemize}
\item Supplied local UnsolvedMath registry record, permanent ID \(30004666\), title ``Normal Schubert Varieties in Full Affine Flag Varieties,'' exact source problem code \texttt{catalogue entry}, in the Oberwolfach Reports Open Problems collection. The recorded proposed year is 2021; no proposer is identified.
\item Supplied report citation: \emph{Algebraic Groups (2022)}, Oberwolfach Reports, pp.~1087--1148, DOI \texttt{10.4171/owr/2021/20}, \url{https://doi.org/10.4171/owr/2021/20}. The literature-source list calls this ``Oberwolfach Report 20/2021.'' These supplied year designations are retained without resolving their relationship.
\item Supplied literature source: Bieker, ``normality in the tamely ramified case,'' \url{https://arxiv.org/abs/2508.04876}.
\item Supplied literature source: Bieker--Richarz, ``normality of affine Schubert varieties,'' \url{https://arxiv.org/abs/2312.11154}.
\item Supplied literature assessment, checked on exactly \texttt{2026-08-22}: status ``partially solved.'' The supplied assessment reports normality criteria and complete classifications in important cases, including adjoint groups in sufficiently large characteristic and semisimple-rank-one examples, while stating that the full affine-flag classification for all listed pairs remains incomplete. This is supplied dated curation derived from the Oberwolfach material, not an independently verified literature assessment; the time-sensitive status and cited claims require independent verification.
\item The registry labels the statement \texttt{corrected\_verified} and records that it was checked against the official Oberwolfach report and preserves the report's hypotheses, notation, and target. This is a supplied source-verification assertion, not a verification performed for this draft.
\end{itemize}
\end{document}
